\documentclass[10pt,a4paper]{article}
\usepackage[a4paper,left=14mm,right=14mm,top=18mm,bottom=15mm,headheight=19pt]{geometry}
\usepackage[T1]{fontenc}
\usepackage[utf8]{inputenc}
\usepackage[scaled=0.94]{helvet}
\renewcommand{\familydefault}{\sfdefault}
\usepackage{xcolor}
\usepackage{tabularx}
\usepackage{booktabs}
\usepackage{array}
\usepackage{fancyhdr}
\usepackage{lastpage}
\usepackage{tcolorbox}
\tcbuselibrary{skins}
\usepackage{hyperref}

% Clinical Color Palette
\definecolor{Navy}{HTML}{16324F}
\definecolor{Teal}{HTML}{0E7490}
\definecolor{SlateBg}{HTML}{F3F6FA}
\definecolor{WarnBg}{HTML}{FFFBEB}
\definecolor{WarnBorder}{HTML}{D97706}
\definecolor{FieldBg}{HTML}{F8FAFC}
\definecolor{BorderGray}{HTML}{B6C2CF}
\definecolor{Muted}{HTML}{475569}
\definecolor{RuleGray}{HTML}{CBD5E1}

\hypersetup{
  colorlinks=true,
  urlcolor=Teal,
  linkcolor=Navy,
  pdftitle={PET Scan Imaging Forms — Fillable Clinical Workbook},
  pdfauthor={Medical Imaging Workbook}
}

% Header & Footer
\pagestyle{fancy}
\fancyhf{}
\fancyhead[L]{\fontsize{8}{10}\selectfont\bfseries\color{Navy} PET SCAN IMAGING FORMS}
\fancyhead[R]{\fontsize{8}{10}\selectfont\color{Muted} Positron Emission Tomography (PET/CT \& PET/MRI) Workflow}
\fancyfoot[L]{\fontsize{7.5}{9}\selectfont\color{Muted} De-identified clinical \& educational use}
\fancyfoot[R]{\fontsize{7.5}{9}\selectfont\color{Muted} Page \thepage\ of \pageref{LastPage}}
\renewcommand{\headrulewidth}{0.5pt}
\renewcommand{\footrulewidth}{0.4pt}
\renewcommand{\headrule}{\hbox to\headwidth{\color{RuleGray}\leaders\hrule height \headrulewidth\hfill}}
\renewcommand{\footrule}{\hbox to\headwidth{\color{RuleGray}\leaders\hrule height \footrulewidth\hfill}}

\setlength{\parindent}{0pt}
\setlength{\parskip}{3pt}

% Interactive Form Field Helpers
\newcommand{\FormLabel}[1]{{\fontsize{7.8}{9.5}\selectfont\bfseries\color{Navy}#1}}

\newcommand{\FText}[3][15pt]{%
  \TextField[name=#2, width=#3, height=#1, bordercolor=BorderGray, backgroundcolor=FieldBg, charsize=8.5pt, borderwidth=0.6pt]{}%
}

\newcommand{\FMText}[3][30pt]{%
  \TextField[name=#2, width=#3, height=#1, multiline=true, bordercolor=BorderGray, backgroundcolor=FieldBg, charsize=8.5pt, borderwidth=0.6pt]{}%
}

\newcommand{\FCheck}[2]{%
  \CheckBox[name=#1, width=9pt, height=9pt, bordercolor=Navy, backgroundcolor=FieldBg, borderwidth=0.7pt]{}~{\fontsize{8}{10}\selectfont #2}%
}

\newcommand{\FieldRow}[2]{%
  \FormLabel{#1}\par\vspace{1pt}%
  \FText[15.5pt]{#2}{\linewidth}\par\vspace{3.5pt}%
}

\newcommand{\MFieldRow}[3][28pt]{%
  \FormLabel{#2}\par\vspace{1pt}%
  \FMText[#1]{#3}{\linewidth}\par\vspace{3.5pt}%
}

\newcommand{\PageBanner}[3]{%
  \noindent\begin{tcolorbox}[
    enhanced,
    colback=Navy,
    colframe=Navy,
    arc=2pt,
    boxrule=0pt,
    left=7pt,right=7pt,top=4.5pt,bottom=4.5pt,
    after skip=4pt
  ]
    {\fontsize{11.5}{14}\selectfont\bfseries\color{white}#1}
    \hfill
    \tcbox[on line, arc=2pt, colback=Teal, colframe=Teal, boxrule=0pt, left=5pt, right=5pt, top=1.5pt, bottom=1.5pt]{%
      \fontsize{7.5}{9}\selectfont\bfseries\color{white}#2%
    }
    \quad
    {\fontsize{8.5}{10}\selectfont\color{white!85!Navy} #3}
  \end{tcolorbox}%
}

\newtcolorbox{GuidanceBox}[1][]{
  enhanced,
  colback=SlateBg,
  colframe=RuleGray,
  arc=1.5pt,
  boxrule=0.5pt,
  left=6pt,right=6pt,top=4pt,bottom=4pt,
  after skip=4pt,
  #1
}

\newcommand{\RefBar}[1]{%
  \vfill
  \noindent\begin{tcolorbox}[
    enhanced,
    colback=SlateBg,
    colframe=RuleGray,
    arc=1.5pt,
    boxrule=0.4pt,
    left=6pt,right=6pt,top=3pt,bottom=3pt,
    after skip=0pt
  ]
    {\fontsize{7.5}{9.5}\selectfont\bfseries\color{Muted} References:\quad #1}
  \end{tcolorbox}%
}

\begin{document}
\begin{Form}

% ============================================================================
% PAGE 1: PET SCAN REFERRAL, PREP & RADIOTRACER SAFETY
% ============================================================================
\PageBanner{PET Scan (PET/CT \& PET/MRI)}{1. BEFORE THE SCAN}{Form 1 of 4}

\begin{GuidanceBox}
{\fontsize{8.1}{10.1}\selectfont
\textbf{\color{Navy}What to check first:} Confirm pathology/biopsy status, clinical indication (staging, restaging, response, infection, neuro/cardiac), prior CT/MRI/PET, and recent chemo/radiation/steroids/G-CSF.\par\vspace{1.5pt}
\textbf{\color{Navy}When to proceed or defer:} Proceed when metabolic/molecular imaging will change management. Defer or reschedule if recent therapy/hyperglycaemia would confound FDG uptake.\par\vspace{1.5pt}
\textbf{\color{Navy}Tracer \& prep safety:} Review pregnancy/breastfeeding and pediatric weight-based dosing. Check fasting (4--6 h for FDG), blood glucose target, and renal function if diagnostic IV contrast CT is planned.}
\end{GuidanceBox}

\noindent\begin{tabularx}{\linewidth}{@{}X@{\hspace{6pt}}p{2.8cm}@{\hspace{6pt}}p{2.4cm}@{\hspace{6pt}}X@{}}
\FormLabel{Patient ID} & \FormLabel{Age / Weight (kg)} & \FormLabel{Date} & \FormLabel{Ordering clinician} \\
\FText[15pt]{pet_p1_pid}{\linewidth} & \FText[15pt]{pet_p1_wt}{\linewidth} & \FText[15pt]{pet_p1_date}{\linewidth} & \FText[15pt]{pet_p1_clin}{\linewidth} \\
\end{tabularx}\par\vspace{3pt}

\noindent\begin{tabularx}{\linewidth}{@{}X@{\hspace{6pt}}p{5.4cm}@{}}
\FormLabel{PET tracer \& scan coverage requested} & \FormLabel{Fasting, glucose, insulin \& recent therapy dates} \\
\FText[15pt]{pet_p1_region}{\linewidth} & \FText[15pt]{pet_p1_prep}{\linewidth} \\
\end{tabularx}\par\vspace{3pt}

\noindent\begin{tcolorbox}[enhanced,colback=SlateBg,colframe=RuleGray,arc=1.5pt,boxrule=0.5pt,left=6pt,right=6pt,top=3.5pt,bottom=3.5pt,after skip=3.5pt]
\FormLabel{Tracer:}\quad
\FCheck{pet_p1_tfdg}{$^{18}$F-FDG}\quad \FCheck{pet_p1_tpsma}{PSMA}\quad \FCheck{pet_p1_tdot}{DOTATATE}\quad \FCheck{pet_p1_tneu}{Amyloid/Tau}\quad \FCheck{pet_p1_toth}{Other}\qquad
\FormLabel{CT:}\quad
\FCheck{pet_p1_ctld}{Low-dose AC}\quad \FCheck{pet_p1_ctiv}{Diagnostic IV contrast}\par\vspace{2.5pt}
\FormLabel{Indication:}\quad
\FCheck{pet_p1_istg}{Initial staging}\quad \FCheck{pet_p1_irstg}{Restaging}\quad \FCheck{pet_p1_iresp}{Response}\quad \FCheck{pet_p1_inod}{Nodule/FUO}\quad \FCheck{pet_p1_inc}{Neuro/Cardiac}\par\vspace{2.5pt}
\FormLabel{Pregnancy:}\quad
\FCheck{pet_p1_pna}{N/A}\quad \FCheck{pet_p1_pneg}{Excluded}\quad \FCheck{pet_p1_ppos}{Pregnant}\quad
\FormLabel{Lactation:}\ \FCheck{pet_p1_lno}{No}\ \FCheck{pet_p1_lyes}{Yes (plan)}\quad
\FormLabel{Pediatric:}\ \FCheck{pet_p1_ped}{Weight-adjusted}
\end{tcolorbox}

\FieldRow{Clinical question to be answered}{pet_p1_cquest}

\MFieldRow[22pt]{History received (exact account, source \& relationship, recipient, date \& time)}{pet_p1_hist}

\MFieldRow[22pt]{Pathology / biopsy confirmation \& prior CT / MRI / PET results (with dates)}{pet_p1_assess}

\noindent\begin{tabularx}{\linewidth}{@{}X@{\hspace{6pt}}X@{}}
\FormLabel{Guideline and exact scenario matched} & \FormLabel{Finding supporting ordering or deferral} \\
\FText[15pt]{pet_p1_guide}{\linewidth} & \FText[15pt]{pet_p1_find}{\linewidth} \\
\end{tabularx}\par\vspace{3pt}

\noindent\begin{tcolorbox}[enhanced,colback=SlateBg,colframe=RuleGray,arc=1.5pt,boxrule=0.5pt,left=6pt,right=6pt,top=3.5pt,bottom=3.5pt,after skip=3.5pt]
\FormLabel{Decision:}\quad
\FCheck{pet_p1_ord}{Order}\quad \FCheck{pet_p1_def}{Defer}\qquad
\FormLabel{Urgency:}\quad
\FCheck{pet_p1_urout}{Routine}\quad \FCheck{pet_p1_uurg}{Urgent}\quad \FCheck{pet_p1_ustat}{Stat / Emergent}\par\vspace{2.5pt}
\FormLabel{What the result would change (or patient-specific exception)}\par\vspace{1pt}
\FText[14.5pt]{pet_p1_change}{\linewidth}
\end{tcolorbox}

\MFieldRow[21pt]{Safety / glucose / breastfeeding decision and plain-language explanation for the patient}{pet_p1_safety}

\noindent\begin{tcolorbox}[enhanced,colback=SlateBg,colframe=RuleGray,arc=1.5pt,boxrule=0.5pt,left=6pt,right=6pt,top=3.5pt,bottom=3.5pt,after skip=2pt]
\FormLabel{Basis for request:}\quad
\FCheck{pet_p1_bclin}{Supported clinical use}\quad
\FCheck{pet_p1_bunc}{Uncertain clinical benefit}\quad
\FCheck{pet_p1_bres}{Research}\quad
\FCheck{pet_p1_bclar}{Awaiting clarification}\par\vspace{2.5pt}
\FormLabel{If research, study reference \& applicable approvals (or clarification note)}\par\vspace{1pt}
\FText[14.5pt]{pet_p1_resref}{\linewidth}
\end{tcolorbox}

\RefBar{\href{https://pmc.ncbi.nlm.nih.gov/articles/12165998/}{FDG PET/CT Staging Reference} \quad|\quad \href{https://snmmi.org/Web/Clinical-Practice/Procedure-Standards/Procedure-Standards}{SNMMI Procedure Standards} \quad|\quad \href{https://pubmed.ncbi.nlm.nih.gov/28937575/}{Pregnancy \& lactation}}

\newpage

% ============================================================================
% PAGE 2: PET SCAN-SPECIFIC REPORT
% ============================================================================
\PageBanner{PET Scan (PET/CT \& PET/MRI)}{2. SCAN-SPECIFIC REPORT}{Form 2 of 4}

\begin{GuidanceBox}
{\fontsize{8.2}{10.2}\selectfont\textbf{\color{Navy}What this report should record:} Record tracer, administered activity (MBq/mCi), injection-to-scan uptake time, blood glucose, scan coverage, CT technique, relevant node groups and organ sites with SUVmax (plus liver/mediastinal reference if applicable), CT correlates, and limits. Distinguish suspicious from indeterminate uptake.}
\end{GuidanceBox}

\noindent\begin{tabularx}{\linewidth}{@{}X@{\hspace{6pt}}p{2.8cm}@{\hspace{6pt}}X@{}}
\FormLabel{Patient ID} & \FormLabel{Exam date} & \FormLabel{Reporting clinician} \\
\FText[15pt]{pet_p2_pid}{\linewidth} & \FText[15pt]{pet_p2_date}{\linewidth} & \FText[15pt]{pet_p2_rep}{\linewidth} \\
\end{tabularx}\par\vspace{3.5pt}

\noindent\begin{tabularx}{\linewidth}{@{}X@{\hspace{6pt}}X@{}}
\FormLabel{Tracer, activity (MBq / mCi), uptake time \& coverage} & \FormLabel{CT technique (low-dose vs. contrast) \& prior comparisons} \\
\FText[15.5pt]{pet_p2_exam}{\linewidth} & \FText[15.5pt]{pet_p2_comp}{\linewidth} \\
\end{tabularx}\par\vspace{3.5pt}

\FieldRow{Clinical indication / question}{pet_p2_ind}

\noindent\begin{tcolorbox}[enhanced,colback=SlateBg,colframe=RuleGray,arc=1.5pt,boxrule=0.5pt,left=6pt,right=6pt,top=3.5pt,bottom=3.5pt,after skip=3.5pt]
\FormLabel{Technical adequacy:}\quad
\FCheck{pet_p2_tadeq}{Adequate}\quad
\FCheck{pet_p2_tsub}{Suboptimal}\quad
\FCheck{pet_p2_trep}{Repeat needed}\qquad
\FormLabel{Blood glucose (mg/dL or mmol/L):}\ \FText[13.5pt]{pet_p2_gluc}{2.8cm}\par\vspace{2.5pt}
\FormLabel{Limitations (motion, extravasation, hyperglycaemia, brown fat, CT misregistration, coverage)}\par\vspace{1pt}
\FText[15pt]{pet_p2_qual}{\linewidth}
\end{tcolorbox}

\MFieldRow[82pt]{Findings \& measurements (nodal groups, organ sites, SUVmax, liver/mediastinal blood pool reference, CT correlates, response score)}{pet_p2_find}

\MFieldRow[50pt]{Impression (direct answer to the clinical question; distinguish suspicious vs. indeterminate uptake)}{pet_p2_imp}

\noindent\begin{tcolorbox}[enhanced,colback=SlateBg,colframe=RuleGray,arc=1.5pt,boxrule=0.5pt,left=6pt,right=6pt,top=3.5pt,bottom=3.5pt,after skip=3.5pt]
\FormLabel{Question status:}\qquad
\FCheck{pet_p2_qans}{Answered}\qquad
\FCheck{pet_p2_qpart}{Partly answered}\qquad
\FCheck{pet_p2_qinc}{Inconclusive}\par\vspace{2.5pt}
\FormLabel{Recommended next step:}\quad
\FCheck{pet_p2_nnone}{None}\quad
\FCheck{pet_p2_npet}{Follow-up PET/CT}\quad
\FCheck{pet_p2_nct}{Dedicated CT}\quad
\FCheck{pet_p2_nmri}{MRI}\quad
\FCheck{pet_p2_nus}{Ultrasound}\quad
\FCheck{pet_p2_nbx}{Biopsy / MDT corr.}
\end{tcolorbox}

\noindent\begin{tabularx}{\linewidth}{@{}X@{\hspace{6pt}}p{4.8cm}@{\hspace{6pt}}p{2.8cm}@{}}
\FormLabel{Follow-up action} & \FormLabel{Responsible clinician or team} & \FormLabel{Due date} \\
\FText[15.5pt]{pet_p2_fu}{\linewidth} & \FText[15.5pt]{pet_p2_resp}{\linewidth} & \FText[15.5pt]{pet_p2_due}{\linewidth} \\
\end{tabularx}\par\vspace{3.5pt}

\noindent\begin{tabularx}{\linewidth}{@{}X@{\hspace{6pt}}p{4.0cm}@{\hspace{6pt}}p{4.2cm}@{}}
\FormLabel{Urgent communication: recipient} & \FormLabel{Date \& time} & \FormLabel{Method} \\
\FText[15.5pt]{pet_p2_uwho}{\linewidth} & \FText[15.5pt]{pet_p2_utime}{\linewidth} & \FText[15.5pt]{pet_p2_umeth}{\linewidth} \\
\end{tabularx}

\RefBar{\href{https://snmmi.org/Web/Clinical-Practice/Procedure-Standards/Procedure-Standards}{SNMMI Procedure Standards} \quad|\quad \href{https://www.rsna.org/practice-tools/data-tools-and-standards/radreport-reporting-templates}{RSNA RadReport Templates}}

\newpage

% ============================================================================
% PAGE 3: WHAT THE IMAGE SHOWS AND MEANS
% ============================================================================
\PageBanner{PET Scan (PET/CT \& PET/MRI)}{3. WHAT THE IMAGE SHOWS \& MEANS}{Form 3 of 4}

\renewcommand{\arraystretch}{1.2}
\noindent\begin{tabularx}{\linewidth}{@{}>{\bfseries\color{Navy}}p{3.0cm} >{\color{Navy}}p{5.4cm} X@{}}
\toprule
\textbf{Context} & \textbf{Key PET Finding} & \textbf{Typical Meaning / Interpretation} \\
\midrule
Nodal Staging & Focal uptake mapped to node above/below diaphragm & Contributes to staging and disease extent \\
Response (Deauville) & Uptake vs. mediastinal blood pool \& liver & Supports complete vs. partial metabolic response \\
Organ / Extranodal & Focal parenchymal or marrow uptake & Requires pattern and anatomic CT/MRI correlation \\
Targeted Tracers & Focal PSMA or DOTATATE avidity & Supports prostate or neuroendocrine tumour sites \\
Benign Mimics & Symmetric brown fat, bowel, healing, infection & Physiologic or inflammatory causes of uptake \\
\bottomrule
\end{tabularx}\par\vspace{4pt}

\begin{GuidanceBox}
{\fontsize{8.1}{10.1}\selectfont
\textbf{\color{Navy}What PET leaves unanswered:} Infection, inflammation, granulomatous disease, recent surgery/radiation, and G-CSF marrow stimulation can also raise FDG uptake. Small lesions ($<$8--10 mm) or low-grade tumours may show low avidity; hyperglycaemia reduces lesion-to-background contrast.\par\vspace{2pt}
\textbf{\color{Navy}Worked interpretation example:} Finding: Focal FDG uptake mapped to a node. Interpretation: may indicate lymphoma or nodal metastasis. Correlate atypical sites clinically.\par\vspace{2pt}
\textbf{\color{Navy}Illustration guide (4:3 plate):} Show paired PET MIP/axial and fused PET/CT schematics with the same nodal or organ site marked.}
\end{GuidanceBox}

\noindent\begin{tabularx}{\linewidth}{@{}X@{\hspace{6pt}}p{4.0cm}@{\hspace{6pt}}p{2.6cm}@{\hspace{6pt}}p{2.0cm}@{}}
\FormLabel{Primary finding} & \FormLabel{Site / Node / Organ} & \FormLabel{Measurement} & \FormLabel{Unit (SUVmax, mm)} \\
\FText[15.5pt]{pet_p3_find}{\linewidth} & \FText[15.5pt]{pet_p3_loc}{\linewidth} & \FText[15.5pt]{pet_p3_meas}{\linewidth} & \FText[15.5pt]{pet_p3_unit}{\linewidth} \\
\end{tabularx}\par\vspace{3.5pt}

\noindent\begin{tabularx}{\linewidth}{@{}X@{\hspace{6pt}}p{4.0cm}@{\hspace{6pt}}p{2.6cm}@{\hspace{6pt}}p{2.0cm}@{}}
\FormLabel{Secondary finding (if any)} & \FormLabel{Site / Node / Organ} & \FormLabel{Measurement} & \FormLabel{Unit (SUVmax, mm)} \\
\FText[15.5pt]{pet_p3_find2}{\linewidth} & \FText[15.5pt]{pet_p3_loc2}{\linewidth} & \FText[15.5pt]{pet_p3_meas2}{\linewidth} & \FText[15.5pt]{pet_p3_unit2}{\linewidth} \\
\end{tabularx}\par\vspace{3.5pt}

\noindent\begin{tcolorbox}[enhanced,colback=SlateBg,colframe=RuleGray,arc=1.5pt,boxrule=0.5pt,left=6pt,right=6pt,top=3.5pt,bottom=3.5pt,after skip=3.5pt]
\FormLabel{Hypothesis impact:}\quad
\FCheck{pet_p3_hsup}{Supports hypothesis}\qquad
\FCheck{pet_p3_hless}{Makes it less likely}\qquad
\FCheck{pet_p3_hunres}{Leaves it unresolved}
\end{tcolorbox}

\MFieldRow[36pt]{Plain-language explanation for patient / care team}{pet_p3_plain}

\noindent\begin{tcolorbox}[enhanced,colback=SlateBg,colframe=RuleGray,arc=1.5pt,boxrule=0.5pt,left=6pt,right=6pt,top=4pt,bottom=4pt,after skip=2pt]
\FormLabel{Audit Challenge}\par\vspace{2pt}
\noindent\begin{tabularx}{\linewidth}{@{}p{5.5cm}@{\hspace{6pt}}X@{}}
\FormLabel{Reviewer \& date} & \FormLabel{Finding disputed} \\
\FText[15pt]{pet_p3_arev}{\linewidth} & \FText[15pt]{pet_p3_adisp}{\linewidth} \\
\end{tabularx}\par\vspace{2.5pt}
\MFieldRow[24pt]{Supporting evidence}{pet_p3_aevid}
\MFieldRow[24pt]{Response or resolution}{pet_p3_ares}
\end{tcolorbox}

\RefBar{\href{https://www.radiologyinfo.org/en/info/pet}{RadiologyInfo: PET/CT} \quad|\quad \href{https://pmc.ncbi.nlm.nih.gov/articles/12165998/}{FDG PET/CT Reference}}

\newpage

% ============================================================================
% PAGE 4: FIND SIMILAR CASES AND CHECK THE EVIDENCE
% ============================================================================
\PageBanner{PET Scan (PET/CT \& PET/MRI)}{4. SIMILAR CASES \& EVIDENCE}{Form 4 of 4}

\begin{GuidanceBox}[colback=WarnBg,colframe=WarnBorder!50!white,borderline west={3pt}{0pt}{WarnBorder}]
{\fontsize{8.1}{10.1}\selectfont\textbf{\color{Navy}Optional review aid:} Use de-identified clinical terms only. Keep patient names, identifiers, reports, and patient images out of public searches. Record examples for review without delaying care.}
\end{GuidanceBox}

\begin{GuidanceBox}
{\fontsize{8.1}{10.2}\selectfont
\textbf{\color{Navy}1. Check the decision using only pre-scan information (example templates):}\par
\texttt{<diagnosis / indication> <staging / response> <FDG / PSMA> PET CT guideline}\par\vspace{3pt}
\textbf{\color{Navy}2. Find a teaching case using the imaging findings (example templates):}\par
\texttt{<diagnosis> <FDG / PSMA> PET CT <nodal / organ uptake pattern> teaching case}}
\end{GuidanceBox}

\FieldRow{Actual search terms used}{pet_p4_terms}

\noindent\begin{tabularx}{\linewidth}{@{}X@{\hspace{6pt}}p{3.4cm}@{}}
\FormLabel{Source link / citation} & \FormLabel{Access date} \\
\FText[15.5pt]{pet_p4_src}{\linewidth} & \FText[15.5pt]{pet_p4_date}{\linewidth} \\
\end{tabularx}\par\vspace{5pt}

\MFieldRow[62pt]{What matches and what differs}{pet_p4_match}

\MFieldRow[62pt]{How the diagnosis was established (or whether confirmation was unclear)}{pet_p4_diag}

\MFieldRow[62pt]{What the source adds and what remains uncertain}{pet_p4_adds}

\RefBar{\href{https://snmmi.org/Web/Clinical-Practice/Procedure-Standards/Procedure-Standards}{SNMMI Procedure Standards} \quad|\quad \href{https://radiopaedia.org/}{Radiopaedia (PET/CT Teaching Cases)}}

\end{Form}
\end{document}
